% Section 5.2: Dark Matter Substructure in Lambda CDM
% Auto-generated from markdown drafts

\section{Dark Matter Substructure in \texorpdfstring{$\Lambda$CDM}{Lambda CDM}}
\label{sec:5.2}

Before searching for dark matter subhalos among gamma-ray sources, we must first understand what $\Lambda$CDM predicts for the subhalo population of a Milky Way-like galaxy: how many subhalos are expected, at what masses, and how their abundance is shaped by baryonic processes.
This section reviews these predictions, establishing the theoretical foundation for the detectability analysis that follows in Section~\ref{sec:5.3}.

\subsection{Hierarchical Structure Formation}
\label{sec:5.2.1}

A defining prediction of $\Lambda$CDM cosmology is that cosmic structure assembles in a bottom-up fashion.
The primordial power spectrum, processed through the radiation-dominated epoch, retains more power on small spatial scales, so that smaller overdensities reach the nonlinear regime and gravitationally collapse before larger ones~\cite{2012AnP...524..507F}.
The resulting bound objects---dark matter halos---then grow through the continuous accretion of diffuse matter and the merging of pre-existing smaller halos, a process known as hierarchical structure formation~\cite{2020NatRP...2...42V}.

When a smaller halo is accreted by a larger host, it does not dissolve entirely into a smooth background.
Because early-forming halos are highly concentrated, their dense inner cores can withstand the tidal forces of the host and survive as distinct, gravitationally bound substructures~\cite{Bullock:2017xww}.
The outer layers of the infalling halo are preferentially stripped away and absorbed into the ambient host, but the tightly bound core persists as a subhalo.
As a result, $\Lambda$CDM predicts that Milky Way-mass halos are populated by a large number of subhalos spanning a wide range of masses.

The quantitative characterization of this subhalo population has been one of the major achievements of cosmological N-body simulations.
Shortly before the turn of the century, numerical resolution reached the level necessary to robustly demonstrate the survival of subhalos within their host halos~\cite{Bullock:2017xww}.
Since then, dedicated high-resolution simulations have mapped the substructure of Milky Way-like galaxies in detail.
Two benchmark projects are the Via Lactea II (VL-II) simulation~\cite{Diemand:2008in}, which resolves subhalos down to approximately $10^4\,M_\odot$, and the Aquarius project~\cite{Springel:2008cc}, which achieves comparable resolution and has been widely used to calibrate subhalo abundance predictions.

These simulations reveal a subhalo mass function that rises steeply toward lower masses, implying that the total number of subhalos in a Milky Way-like galaxy is vast~\cite{Bullock:2017xww,Cirelli:2024ssz}.
The theoretical mass function extends far below the resolution floor of current simulations; for a canonical WIMP, the smallest bound structures could have masses as low as $\sim 10^{-6}\,M_\odot$~\cite{Cirelli:2024ssz}.
The quantitative predictions for the subhalo population, including the J-factor distributions used in this chapter, are derived from the VL-II simulation extended through repopulation techniques~\cite{2024MNRAS.530.2496A} and are discussed in detail in Section~\ref{sec:sim-jfactor}.

\subsection{Luminous Satellites vs.\ Dark Subhalos}
\label{sec:5.2.2}

Not all subhalos are alike from an observational standpoint.
Above a mass threshold of approximately $10^8\,M_\odot$, a subhalo possesses sufficient gravitational potential to accrete and retain baryonic matter, enabling star formation and producing a visible dwarf galaxy~\cite{Zavala:2019gpq}.
Roughly 50 such satellite galaxies are currently known within the virial radius of the Milky Way~\cite{Bullock:2017xww}, ranging from the classical dwarf spheroidals---Draco, Sculptor, Fornax---to the ultra-faint systems discovered over the past two decades by digital sky surveys such as the Sloan Digital Sky Survey and the Dark Energy Survey~\cite{Bullock:2017xww}.
These dwarf spheroidal galaxies are among the most dark-matter-dominated objects known, with mass-to-light ratios exceeding $\sim 1000$ in some cases, and they have been extensively used as targets for indirect dark matter searches in gamma rays (cf.
\ Section~\ref{sec:targets}).

Below $\sim 10^8\,M_\odot$, however, galaxy formation becomes progressively inefficient.
The ultraviolet background produced by cosmic reionization suppresses gas accretion in halos with virial masses below $\sim 10^9\,M_\odot$, while atomic cooling ceases to operate below $\sim 10^8\,M_\odot$~\cite{Bullock:2017xww}.
The vast majority of dark matter subhalos therefore contain no stars, no gas, and produce no conventional electromagnetic signature.
These are the truly ``dark'' subhalos: structures whose existence is predicted with high confidence by $\Lambda$CDM but which cannot be detected through traditional astronomical observations.
If dark matter annihilates, however, these objects could shine in gamma rays, and it is precisely this population that we target in the analysis presented in this chapter.

The apparent discrepancy between the predicted abundance of subhalos and the observed number of satellite galaxies was historically known as the ``missing satellite problem''~\cite{Bullock:2017xww}.
This tension has been largely resolved on two fronts.
Observationally, improved surveys have steadily increased the census of ultra-faint dwarf galaxies, and statistical corrections for sky coverage and detection efficiency suggest that the total satellite population is larger than the pre-2004 sample indicated~\cite{Bullock:2017xww}.
Theoretically, hydrodynamical simulations including baryonic physics---UV feedback, supernova-driven outflows, reionization---naturally suppress galaxy formation in low-mass halos, bringing predictions into broad agreement with observations once detection efficiencies are accounted for~\cite{Cirelli:2024ssz}.

Baryonic processes also modify the survival of the subhalos themselves.
Tidal stripping and gravitational shocks from the Galactic disk preferentially deplete the inner Galaxy of substructure~\cite{2020MNRAS.491.4591E,2018MNRAS.474.3043V,Cirelli:2024ssz}, and the survival rate of low-mass subhalos below $\sim 10^6\,M_\odot$ remains uncertain~\cite{2023MNRAS.518...93A}.
This uncertainty propagates directly into predictions for the number of gamma-ray-detectable subhalos and is discussed further in Section~\ref{sec:sim-jfactor}.

\begin{figure}[t]
	\centering
	% TODO: insert a figure showing the dark matter substructure around a MW-type galaxy
	% Candidate: Diemand et al. (2008), arXiv:0805.1244, Figure 1 (VL-II)
	% Alternative: Bullock & Boylan-Kolchin (2017), arXiv:1707.04256, Figure 7
	\includegraphics[width=\columnwidth]{example-image-a}
	\caption{Dark matter substructure in a Milky Way-like halo from an N-body cosmological simulation, illustrating the vast population of subhalos predicted by $\Lambda$CDM.
		Subhalos above $\sim 10^8\,M_\odot$ can host dwarf satellite galaxies, while the far more numerous low-mass subhalos remain dark.
		From~\cite{Diemand:2008in}.
	}
	\label{fig:nbody_subhalos}
\end{figure}
